
% AI 工具使用声明
% 依据《全国大学生数学建模竞赛人工智能工具使用规定》(2026 试行) 第 3 条
% 必须放在 9.参考文献.tex 之前
\section*{AI 工具使用声明}

\begin{quote}
本参赛队在竞赛过程中使用了 AI 工具,主要用于 \textbf{【简要用途:语言润色、代码调试、模型分析、论文排版】},详细使用情况见支撑材料中的 ``AI 工具使用详情.pdf''。
\end{quote}

\vspace{0.5em}

\noindent\textbf{具体使用情况概述:}

\begin{itemize}
  \item \textbf{AI 工具名称:} 【请填写所用 AI 工具名称,如:大语言模型 X】
  \item \textbf{版本或型号:} 【请填写版本号】
  \item \textbf{使用目的:} 【简要描述,如辅助完成问题分析、Python 求解代码编写、LaTeX 论文排版等】
  \item \textbf{使用环节:} 问题分析 $\rightarrow$ 模型比选 $\rightarrow$ Python 求解 $\rightarrow$ LaTeX 论文写作 $\rightarrow$ 验收修改
  \item \textbf{人工审查与核验:} 全部 AI 生成的代码均在本地 Python 环境下运行验证,生成的数学模型、公式推导和结论均经人工复核;AI 生成的 LaTeX 文本经过逐字人工校对与修改,所有数据结果与原始代码运行结果一致
\end{itemize}

\noindent\textbf{关于核心建模的声明:} 本论文的核心建模思路、模型选择、参数设置、结果分析均由参赛队主导完成,AI 工具仅作为辅助手段,所有最终结论由参赛队独立负责。

\noindent\textbf{关于参赛队责任的声明:} 依据《全国大学生数学建模竞赛参赛规则》(2026 修订稿) 第 6 条,本参赛队对所提交作品的**原创性、真实性和准确性负全部责任**。若论文中存在任何由 AI 工具辅助生成但未充分核验的内容,或与 AI 实际参与情况不符的描述,均由本参赛队承担相应责任。

\vspace{1em}

\noindent 详细使用情况记录在支撑材料中的 ``AI 工具使用详情.pdf'' 文件中,内容应包括:
\begin{enumerate}
  \item 所用 AI 工具名称、版本或型号
  \item 具体使用目的和环节
  \item 主要提示方式与使用过程说明(含典型交互示例)
  \item 对 AI 输出的采纳、人工修改和核验情况
\end{enumerate}
